% Imporditud: tools/import_eksperimendid.py
% Allikas: Eesti füüsikaolümpiaadi eksperimendi ülesannete
%   kogu (Rael Kalda, 2023), ülesanne Ü9, lahendus L9.
\setAuthor{EFO žürii}
\setRound{piirkonnavoor}
\setYear{2002}
\setNumber{P E2}
\setDifficulty{1}
\setTopic{dünaamika}
\setVahendid{kumminiit, koormis massiga 100 g, joonlaud, statiiv}
\setPunktid{12}
\setAllikas{Eksperimendi kogu Ü9, lahendus lk 37}
\setLahendusseis{toores}

\prob{Kumminiit}
Määrata kui suur töö tehakse kumminiidi venitamisel jõuga 1 N.

\hint

\solu
Kinnitada kumminiit statiivi klambri külge. Mõõta kumminiidi esialgne pikkus. Asetada kumminiidi otsa koormis ja mõõta kumminiidi pikkus. Arvutada kumminiidi pikenemine ehk koormise langemise kõrgus. Katset korrata mitu korda ja arvutada keskmine pikenemine. Arvutada töö, mille tegi raskusjõud koormise liikumisel allapoole. Teisendada ühikud. Vastavalt energia jäävuse seadusele on raskusjõu töö võrdne kumminiidi venitamisel tehtava tööga.

Hindamine: Katse kirjeldus --- 3 p. Kumminiidi esialgse pikkuse mõõtmine --- 1 p. Koormatud kumminiidi pikkuse mõõtmine --- 1 p. Kumminiidi pikenemise määramine --- 1 p. Kordusmõõtmised --- 1 p. Aritmeetilise keskmise arvutamine --- 1 p. Töö arvutamine --- 1 p. Ühikute teisendamine --- 1 p. Energia jäävuse seaduse rakendamine --- 1 p. Tõepärane tulemus --- 1 p.
\probend
